\documentclass[11pt,a4paper]{article}
\usepackage[utf8]{inputenc}
\usepackage[T1]{fontenc}
\usepackage[margin=2.5cm]{geometry}
\usepackage{amsmath,amssymb}
\usepackage{booktabs}
\usepackage{array}
\usepackage{hyperref}

\newcommand{\tx}[1]{\texttt{#1}}

\title{NSUCRYPTO 2026, Round 2\\Problem 11: ``Who is here?''}
\date{}
\author{Jacob Lu, Raymond Ouyang, Ruiyang Ye}
\begin{document}
\maketitle

\section*{Answer}

\begin{quote}
The 17 labels scattered over the picture, reassembled, spell
\begin{center}
\Large \tx{SALTYKOV-SHCHEDRIN 200!}
\end{center}
\end{quote}

\section{Reassembling the message}

This is an anagram embedded in the picture, carrying 17 small labels placed on specific objects of a village
scene. In row order, left to right, they read

\[
\tx{SH}\quad \tx{RI}\quad \tx{CH}\quad \tx{S}\quad \tx{K}\quad \tx{O}\quad
\tx{2}\quad \tx{T}\quad \tx{0}\quad \tx{L}\quad \tx{ED}\quad \tx{N}\quad
\tx{!}\quad \tx{0}\quad \tx{A}\quad \tx{Y}\quad \tx{V-}.
\]

The combination match
\tx{SALTYKOV-SHCHEDRIN 200!} exactly --- 22 characters either way, with nothing
missing and nothing left over. 

\section{Who is here}

Saltykov-Shchedrin (1826--1889) is here. The picture is a crop of Rotraut
Susanne Berner's \emph{Sommer-Wimmelbuch}, with the 17 labels added on top.

\end{document}